\chapter{Peripheral Angioplasty and Stenting}

\section*{Introduction \& Clinical Indications}
Peripheral angioplasty uses catheters and balloons to improve blood flow through a narrowed or blocked peripheral artery. A stent may be added when balloon treatment alone is inadequate or when vessel recoil or dissection threatens patency. Revascularization is considered for lifestyle-limiting claudication despite medical therapy and structured exercise, or for chronic limb-threatening ischemia when it can relieve ischemia and support wound healing.

Acute limb ischemia is an emergency requiring immediate vascular assessment. Treatment may include thrombectomy, thrombolysis, open bypass, or endovascular intervention. Revascularization should be integrated with smoking cessation, lipid lowering, diabetes and blood-pressure management, antithrombotic therapy, foot protection, and wound care.

\section*{Relevant Surgical Anatomy}
Procedures may involve the aortoiliac, femoropopliteal, or infrapopliteal circulation. Inflow, lesion length, calcification, chronic occlusion, distal runoff, collateral vessels, and the relationship to branch arteries determine strategy. Femoral, radial, or other arterial access routes are selected according to the target and patient anatomy.

\section*{Preoperative Workup \& Preparation}
Document pulses, skin, wounds, rest pain, walking limitation, ankle-brachial or toe pressures, duplex or cross-sectional imaging, renal function, diabetes, smoking, and bleeding risk. Review prior interventions and antiplatelet or anticoagulant therapy. Explain access-site bleeding, pseudoaneurysm, dissection, perforation, distal embolization, contrast-associated kidney injury, acute occlusion, restenosis, limb ischemia, and possible conversion to open treatment.

\section*{Surgical Technique \& Procedure}
After local anesthesia and monitoring, arterial access is obtained and a guidewire is advanced across the lesion under imaging. Balloon angioplasty dilates the stenosis; selected lesions may receive a stent or drug-coated technology according to anatomy, evidence, and local practice. Completion angiography checks flow, branch vessels, dissection, perforation, and distal embolization. Hemostasis is achieved with compression or a closure device when appropriate.

\section*{Complications \& Risks}
Risks include hematoma, pseudoaneurysm, arterial dissection, perforation, distal embolization, acute thrombosis, contrast reaction, kidney injury, restenosis, infection, and worsening limb ischemia. New severe limb pain, pallor, coolness, loss of pulses, neurologic symptoms, chest pain, or uncontrolled bleeding requires emergency assessment.

\section*{Postoperative Care \& Recovery}
Check access sites, distal circulation, renal function, pain, and walking status. Continue prescribed antiplatelet or antithrombotic therapy and risk-factor treatment. Encourage supervised exercise, smoking cessation, foot inspection, wound care, and prompt treatment of infection. Surveillance is guided by symptoms, wound progress, and the treated arterial segment.

\section*{Outcomes \& Prognosis}
Technical success can improve walking capacity, relieve ischemic pain, and support wound healing. Durability varies with lesion anatomy, diabetes, smoking, renal disease, runoff, and adherence to medical therapy. Restenosis and progression remain possible, so revascularization does not replace long-term vascular prevention.

\section*{References}
\begin{enumerate}
\item American Heart Association/American College of Cardiology. Guideline for the Management of Patients With Lower Extremity Peripheral Artery Disease.
\item European Society for Vascular Surgery. Guidelines on chronic limb-threatening ischemia and peripheral arterial disease.
\item Society for Vascular Surgery. Practice guidance for lower-extremity arterial disease.
\end{enumerate}
